\documentclass[manuscript,review,anonymous]{acmart}

%% ------------------------------------------------------------------
%% BASiS 日本語版ドラフト（英語版の内容確認用）
%% pdfLaTeX + CJKutf8 でコンパイルする日本語確認版です。
%% コンパイル順: pdflatex -> bibtex -> pdflatex -> pdflatex
%% ------------------------------------------------------------------

\usepackage{amsmath}
\usepackage{booktabs}
\usepackage{graphicx}

%% Japanese support compatible with pdfLaTeX/Overleaf free-plan builds.
%% CJKutf8 uses TeX Live's bundled Japanese fonts, so no system-font search is needed.
\usepackage{CJKutf8}
\usepackage{CJKspace}

%% Load the bundled Mincho font definition early so acmart's bold/italic
%% requests can be mapped to shapes available under pdfLaTeX.
\makeatletter
\input{c70min.fd}
\makeatother
\DeclareFontShape{C70}{min}{b}{n}{<->ssub*min/bx/n}{}
\DeclareFontShape{C70}{min}{m}{it}{<->ssub*min/m/n}{}
\DeclareFontShape{C70}{min}{b}{it}{<->ssub*min/bx/n}{}

%% Activate CJK decoding for the whole document without opening a nested
%% environment. This keeps acmart's title/keyword metadata assignments global
%% and also keeps CJK active while deferred floats and end-document hooks run.
\makeatletter
\newcommand{\ActivateCJKForPDFLaTeX}{%
  \CJKnospace
  \CJK@envStart{}{UTF8}{min}%
}
\makeatother

\setcopyright{acmlicensed}
\acmConference[IUI '27]{ACM International Conference on Intelligent User Interfaces}{March 2027}{Zurich, Switzerland}
\acmDOI{}
\acmISBN{}

\begin{document}
\ActivateCJKForPDFLaTeX

\title{BASiS: 小規模高品質コーパスの会話スタイルを目的特化対話システムへ反映するためのベイズ的データ選別（日本語版・内容確認用ドラフト）}

\author{Anonymous Author(s)}
\affiliation{%
  \institution{Institution withheld for review}
  \city{City}
  \country{Country}
}
\renewcommand{\shortauthors}{Anonymous Author(s)}

\begin{abstract}
大規模言語モデル（LLM）は流暢な対話と正確な応答を実現できるようになったが，感情支援・チュータリング・医療問診のような目的特化対話では，成否は「何を答えるか」だけでなく「どう応答するか」という会話スタイルに大きく依存する。こうしたスタイルは小規模で高品質なコーパスに最もよく表れるが，そのままでは学習に用いるには小さすぎる。一方，学習に十分な量を持つ大規模・汎用コーパスは，目的の会話スタイルに合致するとは限らない。本研究では，BASiS（Bayesian Style-based Selection）を提案する。BASiSは，(i) LLMを用いて小規模高品質コーパスを会話状態・応答戦略・状態遷移として構造化し，(ii) その構造をベイズ的対話モデル（状態遷移確率と観測尤度）として表現し，(iii) ベイズ更新によって大規模汎用コーパス中の任意の応答候補を目的スタイルとの適合度でスコアリングし，(iv) 得られたスコアに基づいてLoRAベースのDirect Preference Optimization（DPO）用のchosen/rejectedペアを構築してローカルLLMを学習する。既存のガイドラインに基づく応答制御や単一戦略予測とは異なり，BASiSは人手によるガイドライン設計を必要とせず，同一状況下で複数の応答戦略が妥当となる自然な多様性を，単一ラベルではなく確率分布として表現する。ESConvをスタイル参照元，DailyDialogを選別対象プールとした予備実験では，Qwen3.5-27BをBASiSで選別したデータで学習したモデルが，ベースモデルおよびランダム選別モデルの双方に対して，ESConv参照の評価軸およびトップカンファレンスの評価実践に基づく5軸評価（Style Strength，ESConv Tone Similarity，Supporter Role Consistency，Non-directive Support Style，Premature Advice Avoidance）の双方で，Oracle LLMによるブラインド評価においてスタイル適合度の向上を示した。Friedman検定とHolm補正付き対応ありpermutation testにより，ほとんどの軸で統計的に有意な差が確認された。さらに，追加のコーパスペアでの実験と人手評価を現在進めており，本論文ではBASiSが現時点で示す「受容と会話進行のトレードオフ」についても率直に議論する。
\end{abstract}

\begin{CCSXML}
<ccs2012>
 <concept>
  <concept_id>10010147.10010178.10010179</concept_id>
  <concept_desc>Computing methodologies~Natural language generation</concept_desc>
  <concept_significance>500</concept_significance>
 </concept>
 <concept>
  <concept_id>10010147.10010178.10010179.10010182</concept_id>
  <concept_desc>Computing methodologies~Discourse, dialogue and pragmatics</concept_desc>
  <concept_significance>500</concept_significance>
 </concept>
 <concept>
  <concept_id>10003120.10003121.10003122.10003334</concept_id>
  <concept_desc>Human-centered computing~Natural language interfaces</concept_desc>
  <concept_significance>300</concept_significance>
 </concept>
 <concept>
  <concept_id>10010147.10010257.10010293.10010294</concept_id>
  <concept_desc>Computing methodologies~Learning from implicit feedback</concept_desc>
  <concept_significance>100</concept_significance>
 </concept>
</ccs2012>
\end{CCSXML}

\ccsdesc[500]{Computing methodologies~Natural language generation}
\ccsdesc[500]{Computing methodologies~Discourse, dialogue and pragmatics}
\ccsdesc[300]{Human-centered computing~Natural language interfaces}
\ccsdesc[100]{Computing methodologies~Learning from implicit feedback}

\keywords{目的特化対話システム, 会話スタイル, データ選別, ベイズ的対話モデリング, 選好最適化, 感情支援対話, 大規模言語モデル}

\maketitle

%% Keep PDF metadata ASCII-only. acmart writes metadata again at end of the
%% document, after the CJK environment has closed; Japanese metadata would
%% otherwise trigger a pdfLaTeX Unicode error at \end{document}.
\hypersetup{
  pdftitle={BASiS: Japanese Content-Review Draft},
  pdfauthor={Anonymous Author(s)},
  pdfsubject={Purpose-specific dialogue style adaptation},
  pdfkeywords={dialogue systems, conversational style, data selection, Bayesian modeling, DPO}
}

\section{はじめに}
\label{sec:intro}

Instruction tuningやRLHFにより，近年の大規模言語モデル（LLM）は非常に高い一般対話能力を獲得した。すなわち，指示に従い，正確に回答し，多様な話題にわたって自然で一貫した複数ターンの対話を維持できる\cite{ouyang2022instructgpt}。しかし，実務上重要な対話アプリケーションの多く――感情支援，カウンセリング的対話，チュータリング，医療問診など――は，主として事実的な正確さや流暢さでは評価されない。それらは応答の\emph{仕方}によって評価される。すなわち，解決策を提示する前にユーザの感情を受け止めているか，複数ターンにわたって一貫した支援的役割を維持できているか，最速の質問ではなく適切なタイミングで適切な質問をしているか，といった点である\cite{rashkin2019empathetic,liu2021esconv}。このような場面では，「何を答えるか」と同じくらい「どう応答するか」が重要であり，一般対話能力がいかに高くとも，それだけでは適切な会話スタイルを保証しない。

LLMを目的の会話スタイルに適応させることは，本質的にはデータの問題であり，利用可能なデータには互いに補完的な弱点を持つ2つの領域が存在する。大規模・汎用の対話コーパスは豊富で多様だが，特定の目的のために収集されたものではないため，目的の会話スタイルに合致する部分はごく一部にとどまり，かつその部分を容易に切り出せる保証もない。一方，小規模で目的特化的に構築されたコーパスは，明示的に目的の会話スタイルを反映するよう収集されており品質は高いが，収集には専門知識を要し，特にカウンセリングや医療のような機微な領域ではプライバシー上の制約もあり，収集コストが高く時間もかかる\cite{zhou2023lima,zhang2025datasurvey}。その結果，小規模高品質コーパスは，そのまま現代のLLMを学習させるには小さすぎることが多く，直接学習すればごく一部の表層的パターンへの過学習を招くか，継続的な収集コストを要することになる。

この緊張関係は，自然な問いを生む。すなわち，小規模高品質コーパスの持つスタイル的特徴を，大規模汎用コーパスからの\emph{データ選別}に活用し，学習データをスタイル適合的かつ十分な量を確保した状態にできないか，という問いである。既存研究には，この問いに部分的に答える2つの系統があるが，それぞれ構造的な限界を抱えている。ガイドラインに基づく応答制御は，望ましい振る舞いを明示的なルールやポリシー事前分布として符号化し，それに従うように生成を再ランキングまたは誘導する\cite{wen2025guidedtod}。この方法は事前に人手でガイドラインを設計する必要があり，設計コストが高いうえ，会話スタイルの微妙な側面――間の取り方，トーン，受容と会話進行のバランスなど――は，まさに明示的なチェックリストとして書き下すのが難しい暗黙知である。一方，応答戦略予測は，対話の現在の状態が与えられたとき，次にどの戦略（例えば共感，質問，助言など）を取るべきかをデータから学習する\cite{wan2025emodynamix}。これは人手によるルール設計を回避できるが，その構造上，次の戦略を単一のラベルとして予測する。実際の目的コーパスには，同じ文脈に対して複数の妥当な戦略が存在することが多く，この多様性を単一の予測ラベルに縮約することは，スタイル適合的なデータ選別手法が本来保持すべき多様性を損なうことになる。

本研究では，これら2つの限界をともに回避するデータ選別手法として\textbf{BASiS}（\textbf{Ba}yesian \textbf{S}tyle-based \textbf{S}election）を提案する。BASiSはまず，LLMを用いて小規模高品質コーパスを読み込み，各ターンを3つの要素――話者の状況・感情・対話段階を表す\emph{会話状態}，共感・質問・助言などの応答者側のアプローチを表す\emph{応答戦略}，そしてその戦略が引き起こす\emph{状態遷移}――として構造化する。この構造化された表現は，状態遷移分布と，応答戦略と状態を関係づける観測尤度からなるベイズ的対話モデルとして符号化される。このモデルを用いて，BASiSは大規模汎用コーパスから抽出した任意の応答候補をスコアリングする。具体的には，これまでの会話履歴を状態上の信念分布として扱い，遷移モデルによって次状態の予測分布を計算し，候補応答から観測される戦略ラベルによってベイズ更新を行い，望ましい目的スタイルに対応する状態群への事後確率質量を合算する。このスコアリング機構は単一の予測ラベルではなく確率分布上で動作するため，ある状況に対して複数の戦略が妥当である事実を自然に表現でき，人手によるガイドライン設計を一切必要としない。得られたスコアはchosen/rejectedの選好ペアの構築に用いられ，LoRAベースのDirect Preference Optimization（DPO）によってローカルLLMを学習することで，フルファインチューニングのコストをかけずに目的スタイルをモデルへ反映する。

本研究では，Emotional Support Conversationコーパス（ESConv）\cite{liu2021esconv}を小規模高品質なスタイル参照元，DailyDialog\cite{li2017dailydialog}を大規模汎用な選別対象プールとした予備実験を行い，2{,}500件のBASiS選別ペアを用いてQwen3.5-27BをLoRAベースDPOで学習した。Oracle LLM（GPT-5.4）による評価では，学習後モデルを，ファインチューニングを行わないベースモデルおよび同一プールからランダムに選別したデータで学習したモデルと比較した。評価は，ESConvデータセット論文自体の評価軸に準拠したルーブリックと，トップカンファレンスの評価実践から構成した5軸ルーブリック――Style Strength，ESConv Tone Similarity，Supporter Role Consistency，Non-directive Support Style，Premature Advice Avoidance\cite{mir2019styletransfer,liu2021esconv,zhang2018persona,deng2023mixedinitiative}――の双方で行い，Friedman検定とHolm補正付き対応ありpermutation testによる有意性検定を併用した。BASiSで選別したデータは，ほぼすべての軸でベースモデル・ランダム選別モデルの双方に対してスタイル適合度を向上させ，とりわけ感情の受容と助言のタイミングに直結する軸で最大の改善を示した。一方で，本手法は現時点でひとつのトレードオフを抱えていることも判明した。すなわち，BASiSは支援的で受容的な振る舞いを強化する一方，会話を積極的に前進させる力については一定のコストを伴う。本論文ではこのトレードオフを率直に報告し，感情支援以外の目的特化領域を対象とした追加実験，および人手評価という進行中の取り組みが，このトレードオフがこの領域固有のものか，より一般的な性質かを検証するためにどのように設計されているかを述べる。

本研究の貢献は次の4点である。(1) 人手によるガイドライン設計を必要とせず，小規模高品質コーパスの暗黙的な会話スタイルをベイズ的対話モデルとして構造化する手法の提案。(2) 複数の妥当な戦略を単一の予測ラベルに縮約せず，確率質量として表現しつつ任意の候補応答をスコアリングするベイズ更新に基づくスコアリング機構の提案。(3) 目的固有および分野横断的なスタイル評価軸と適切な統計的検定を用いた実証実験により，BASiSで選別したデータがベースモデルおよびランダム選別ベースラインの双方に対して会話スタイル適合度を向上させることを示したこと。(4) 本手法の現時点での限界を率直に述べ，感情支援以外の目的特化領域および人手評価に対する妥当性を検証するための具体的かつ既に進行中の計画を提示したこと。

以下，第\ref{sec:relwork}節では，目的特化対話，ガイドラインおよび戦略に基づく応答制御，LLM適応のためのデータ選別，ベイズ的対話モデリング，スタイル評価に関する関連研究をレビューし，BASiSが取り組む具体的なギャップを整理する。第\ref{sec:method}節ではBASiSのパイプラインを詳述する。第\ref{sec:setup}節では実験設定を，第\ref{sec:results}節では実験結果を報告する。第\ref{sec:discussion}節で結果を考察し，第\ref{sec:limitations}節で限界と進行中の取り組みを述べたのち，倫理声明，GenAI利用開示，および結論を述べる。


\section{関連研究}
\label{sec:relwork}

\subsection{目的特化対話と会話スタイル}
Instruction tuningとRLHFにより，一般的なLLMは指示追従性が高く，正確に回答し，流暢な複数ターンの対話を維持できるようになった\cite{ouyang2022instructgpt}。これとは別に，対話システムがある特定の目的に資するとき，この一般的な能力に加えて何が必要かを問う研究の系統がある。共感的なオープンドメイン対話に関する研究は，話者の感情状態を認識し適切に応答する能力が，明示的に狙って獲得させるべき能力であることを示し，それを測定するためのベンチマークデータを導入した\cite{rashkin2019empathetic}。Emotional Support Conversationタスクと付随するESConvコーパス\cite{liu2021esconv}は，これをさらに先鋭化させ，質問・言い換え・感情の反映・提案の提供といった支援戦略の明示的な一覧を定義し，人間の支援者が特徴的な戦略の系列に従うこと――典型的には，解決策を提示する\emph{代わりに}ではなく，その\emph{前に}来訪者の感情を探索し受容すること――を示した。目的特化対話は感情支援に限らない。数学チュータリング対話では，答えをそのまま教えるのではなく，生徒を解法へと足場かけする必要があり\cite{macina2023mathdial}，医療問診対話では，患者が有能で安心できると感じる形で症状の属性を体系的に聞き出す必要がある\cite{saley2024meditod}。これらの互いに大きく異なる領域に共通するのは，成否が内容の正しさの上に重なった\emph{スタイルないし戦略の層}に依存しており，この層は一般対話能力だけでは十分に規定されないという点である。BASiSは，スタイル参照元として用いる領域に依存しない設計を志向しており，第\ref{sec:setup}節では，感情支援からチュータリング・医療問診の組み合わせ，およびより雑音の多い大規模な一般対話データへと評価を拡張している現状を述べ，この領域非依存性を検証する。

\subsection{ガイドラインおよびポリシーに基づく応答制御}
目的の会話スタイルを実現する1つのアプローチは，望ましい振る舞いをガイドラインやポリシーとして明示化し，それに条件付けて生成を行うことである。Wenらは，タスク指向対話の流れを確率的モデルで表現し，そのモデルを用いて従うべきガイドライン準拠のポリシーを逐次更新し，更新後のポリシーへの準拠度に応じて応答候補を再ランキングするGuidedTODを提案した\cite{wen2025guidedtod}。既存のガイドライン準拠手法と比較して，GuidedTODはより高いガイドライン準拠性能を達成している。この結果は，目的特化対話においてモデルの無制御な生成に頼るのではなく，対話の流れを明示的にモデル化することの有効性を示しており，この前提はBASiSも共有するところである。しかし，この手法はシステム構築に先立って人手でガイドラインを設計することを前提としており，これには高いコストがかかる。さらに，会話スタイルを定義する特徴の多くは，明示的というより暗黙的である。ESConvらしい支援者応答が適切な間合いや適切な温かさを感じさせる理由は，やってよいこと・悪いことのチェックリストとして自然に表現できるものではなく，そのようなガイドライン文書は必然的に不完全となり，ガイドライン準拠の生成だけでは埋められないスタイル上の空白を残す。

\subsection{状態・戦略を考慮した応答予測}
第2のアプローチは，人手のガイドラインからではなく，対話そのものから適切な次の応答戦略を推論するものである。Wanらは，ユーザの細かな感情状態とシステムの応答戦略の関係を異種グラフとしてモデル化し，そのグラフを用いて次に取るべき戦略を予測するEmoDynamiXを提案した\cite{wan2025emodynamix}。EmoDynamiXは高い戦略予測性能を達成し，グラフ構造が解釈可能であるためにどの戦略がなぜ選ばれたかについて透明性も提供する。これはガイドライン設計と比較して魅力的な特徴である。しかし，原論文自身も認めているように，中心的な限界は，各ターンにつき単一の次戦略を予測する点にある。実際には，同一の会話状態が複数の適切な戦略を許容することが多い――来訪者の感情の受容を続けることと，明確化のための質問をすることは，同じ文脈でしばしばどちらも妥当な選択となる――ため，この多様性を単一の予測ラベルに縮約することは，目的コーパスが実際に含んでいる情報を失わせる。このような予測器を，たとえ間接的にであれ適応用の候補データのフィルタリングやランキングに用いると，この脆弱性により選別データが1つの支配的な戦略に偏り，参照コーパスに存在する正当なスタイルの多様性から逸れてしまうおそれがある。

\subsection{LLM適応のためのデータ選別}
対話固有の構造とはおおむね独立に，LLM適応のためのデータの選別・フィルタリング方法を問う，別だが関連する研究群も存在する。LIMAは，注意深く選別された指示応答ペアが（おおよそ千件程度という）少数であっても，事前学習済みLLMの振る舞いを整合させるのに十分であり得ることを示し，アライメントの多くは新しい知識の注入ではなく，表層的な形式やスタイルに関するものであると論じている\cite{zhou2023lima}。この知見は，小規模高品質コーパスを\emph{スタイル}シグナルの正当かつ十分な情報源として扱うという，本研究が前提とする考え方を裏付けるものであるが，LIMA自体は選別された小規模データ集合上で直接学習を行うのみであり，そのシグナルを大規模プールへとスケールさせる仕組みは提供していない。AlpaGasusは補完的な問題――大規模だが雑音の多い指示チューニングデータが与えられたとき，強力なLLMを用いて各事例にスカラーの品質スコアを付与し，低スコアの事例を除去する――に取り組み，LLMでフィルタリングした小規模な部分集合が，フィルタリングしていない全データでの学習を上回り得ることを示した\cite{chen2023alpagasus}。これは，LLMに基づく自動スコアリングが大規模な人手によるデータキュレーションの代替として有効であることを示すものであり，BASiSもこの仕組みに依拠する。しかし，AlpaGasusにおける「品質」とは，指示応答の正しさと有用性に関する汎用的かつ単一ターンの判断であり，複数ターンの対話構造をモデル化することも，コーパス固有の会話スタイル分布を対象とすることもない。この2点はいずれも，BASiSのベイズ的対話モデルが捉えることを意図している要素である。指示チューニングのためのデータ選別手法に関するより広範なサーベイも，既存手法の大半が，小規模参照コーパスのスタイル分布への適合ではなく，一般的な応答品質・多様性・タスク網羅性を最適化していることを確認している\cite{zhang2025datasurvey}。選別の後段では，Direct Preference Optimization（DPO）が，強化学習ベースの選好アライメントを，chosen/rejected応答ペア上の直接的かつ閉形式の分類目的関数として再定式化し\cite{rafailov2023dpo}，Low-Rank Adaptation（LoRA）が，事前学習済みの重みを凍結したまま少数の低ランクパラメータのみを学習することで大規模モデルのファインチューニングを扱いやすくしている\cite{hu2021lora}。BASiS自体は，選別したデータをどのアライメントアルゴリズムで消費するかについて特定の手法に依存しないが，本研究ではLoRAベースのDPOを採用する。これは，BASiSが各候補応答に付与するScoreが，同一プロンプトに対する高Score応答と低Score応答というchosen/rejectedペアを自然に生成すること，またLoRAがベースモデルの一般的な能力を損なわない程度に適応を軽量に保てることによる。

\subsection{対話のベイズ的・確率的モデリング}
対話を潜在状態上の確率分布として表現し，新たな証拠が得られるたびにその分布を更新するという発想は，タスク指向の音声対話システムにおいて長い歴史を持つ。Partially Observable Markov Decision Process（POMDP）に基づく対話管理は，真の（観測不能な）対話状態上の確率分布である信念状態を保持し，新たな雑音を含む観測（典型的には音声認識仮説）が得られるたびにベイズ則で更新することで，音声認識誤りに対する頑健性を実現する\cite{young2013pomdp}。ThomsonとYoungは，このフレームワーク内で対話状態の扱いやすいベイズ更新を定式化し，近似的な信念伝播によって，リアルタイムのベイズ的状態追跡が対話管理において実用的になることを示した\cite{thomson2010bayesian}。BASiSは，この研究系統の中核的な仕組み――状態上の信念分布，遷移モデル，観測からのベイズ更新――を借用しつつも，異なる用途に転用している。POMDPに基づく対話管理器が信念状態を用いてシステム自身の次の行動を選択し，不確実性を音声認識や言語理解の雑音に由来するものとして扱うのに対し，BASiSは信念状態を用いて，他者（あるいは別のシステム）によって書かれた\emph{候補応答}を，目的スタイルの特徴的なダイナミクスとの整合性という観点でスコアリングし，不確実性を目的スタイルの話者が次に取り得る行動の自然な多様性に由来するものとして扱う。我々の知る限り，ベイズ的信念更新を対話の\emph{制御}からコーパスレベルのスタイル整合性\emph{スコアリング}（データ選別のための）へと転用した研究はこれまでにない。

\subsection{会話スタイルの評価}
第\ref{sec:setup}節で採用する5軸のスタイルルーブリックは，単一の情報源ではなく，複数の異なる研究領域で developされた評価方法論を土台としている。Mirらは，転移強度・内容保存・流暢性といった軸に沿ってテキストスタイル転移の評価を定式化し，スタイルは表層的な語彙の重なりから推測するのではなく，専用の分類器で測定すべきだと論じている\cite{mir2019styletransfer}。本研究はこの「転移強度／スタイルの強さ」という概念を，応答が目的コーパスのスタイルをどれだけ強く示しているかの測定に適用する。ESConv論文自体は，感情支援に特化した戦略・トーンレベルの基準――戦略選択の適切さ，温かさ，感情の受容に対する助言のタイミング――を定義しており，本研究はこれをトーン類似度と助言タイミングの軸に適用する\cite{liu2021esconv}。PersonaChatコーパスとともに導入された対話エージェントのペルソナ一貫性評価は，システムが単一ターン内だけでなく複数ターンの対話全体を通じて一貫した人格や役割を維持しているかを測定する\cite{zhang2018persona}。本研究はこれを役割一貫性の軸に適用する。最後に，mixed-initiativeな感情支援対話に関する研究は，システムのイニシアチブの取り方（例えば助言への移行）が，来訪者の感情開示に対して反射的に指示的であるのではなく，適切なタイミングで行われているかを明示的に測定しており\cite{deng2023mixedinitiative}，本研究はこれを非指示的支援の軸に適用する。これらの評価実践は個別には，（主にESConvのような）特定タスク向けに構築されたシステムを評価するために設計されたものである。我々の知る限り，これらが，特定の目的スタイルに紐づかない汎用的な\emph{スタイル選別手法}を監査するための単一の多軸プロトコルとしてまとめられたことはこれまでになく，本研究の評価におけるこのルーブリックの役割はまさにその点にある。

\subsection{BASiSが取り組むギャップの整理}
以上の関連研究を総合すると，BASiSの設計を動機づける4つのギャップが浮かび上がる。第一に，ガイドラインに基づく応答制御は，しばしば暗黙的な振る舞いに対して，コストが高く不完全な人手ルールを必要とする\cite{wen2025guidedtod}。BASiSは代わりに，LLMを用いて生の対話テキストからスタイルを自動的に構造化する。第二に，単一戦略の応答予測は，目的スタイル内に存在する妥当な戦略の自然な一対多の多様性を過小に表現する\cite{wan2025emodynamix}。BASiSは代わりに，この多様性を信念分布と連続的なScoreとして表現し，単一の予測ラベルへと縮約することがない。第三に，既存のLLM適応向け自動データ選別手法は，複数ターンかつコーパス固有のスタイルではなく，単一ターンかつ汎用的な品質概念に基づいて動作する\cite{chen2023alpagasus,zhang2025datasurvey}。BASiSのベイズ的対話モデルは，目的コーパスから抽出した複数ターンの状態遷移を明示的に中心に据えて構築されている。第四に，既存のスタイル評価実践は，概ね特定タスク向けであり，1つの目的スタイル向けに構築されたシステムを評価するものであって，領域非依存の選別手法を監査するものではない\cite{mir2019styletransfer,liu2021esconv,zhang2018persona,deng2023mixedinitiative}。本研究はこの目的のための分野横断的なルーブリックを構築し，第\ref{sec:limitations}節で述べるように，追加の領域にわたってこれを実際に検証しつつある。


\section{提案手法：BASiS}
\label{sec:method}

\subsection{全体像}
BASiSは，小規模高品質コーパスの会話スタイルを再利用可能なスコアリング関数へと変換し，そのスコアリング関数を用いて大規模汎用コーパスからスタイル適合的なデータを選別し，ローカルLLMのファインチューニングに用いる。パイプラインは図\ref{fig:pipeline}に示す4つのステップからなる：(1) LLMを用いて小規模コーパスから会話特徴を抽出する，(2) その特徴をベイズ的対話モデルとして符号化する，(3) そのモデルの下でのベイズ的信念更新により，大規模コーパスの応答候補をスコアリングする，(4) 高スコア・低スコアのデータを選別してDPO選好ペアを構築し，LoRAによりローカルLLMを学習する。以下，各ステップを順に説明する。

\begin{figure}[t]
  \centering
  %% \includegraphics[width=\linewidth]{figures/fig1_pipeline_overview}
  \fbox{\parbox{0.9\linewidth}{\centering\vspace{2em}図プレースホルダー：BASiSパイプライン全体像\\（小規模コーパス $\rightarrow$ 特徴抽出 $\rightarrow$ ベイズ的対話モデル $\rightarrow$ 大規模コーパスのスコアリング $\rightarrow$ データ選別 $\rightarrow$ LoRA/DPO学習）\\\texttt{fig1\_pipeline\_overview} をここに挿入してください。\vspace{2em}}}
  \caption{BASiSパイプラインの全体像。小規模高品質コーパスは会話特徴（状態・戦略・遷移）として構造化され，ベイズ的対話モデルとして符号化される。このモデルは大規模汎用コーパスの応答候補を目的スタイルとの整合性に基づいてスコアリングし，最高スコアと最低スコアの候補がLoRAベースDPOファインチューニング用のchosen/rejectedペアとして用いられる。}
  \Description{小規模高品質コーパスと大規模汎用コーパスの2つの入力が，会話特徴分析・ベイズ的対話モデリング・スコアリング・データ選別の4段階のプロセスに流れ込み，最終的にローカルLLMのLoRA/DPOファインチューニングに至るパイプライン図。}
  \label{fig:pipeline}
\end{figure}

\subsection{ステップ1：小規模コーパスからの会話特徴抽出}
BASiSは，小規模高品質コーパスを非構造化テキストとして扱うのではなく，LLMを用いて各対話を読み込み，互いに補完的な3つの次元でラベル付けする（図\ref{fig:features}）。\emph{会話状態}$s$は，対話相手の状況・感情・対話段階（例えば，気持ちを開示している段階，その気持ちを整理している段階，解決策を検討している段階）を捉える。\emph{応答戦略}$a$は，与えられた状態に対して応答者が取るアプローチ（例えば，共感・受容，明確化のための質問，情報提供，助言の提供）を捉える。\emph{状態遷移}は，ある状態に対してどの戦略が適用されたときに，どの後続状態へ至るかを記録する。この抽出手続きを小規模コーパス全体に適用することで，(状態，戦略，次状態)の3つ組の集合が得られ，コーパスの暗黙的な会話スタイルが，人手でルールを列挙することなく，構造化されたデータとして明示化される。この構造化ステップは，パイプライン全体の中で小規模コーパスが直接参照される唯一の箇所であり，以降のすべてのステップは，このステップが生成した構造の上で動作する。

\begin{figure}[t]
  \centering
  %% \includegraphics[width=\linewidth]{figures/fig2_feature_extraction}
  \fbox{\parbox{0.9\linewidth}{\centering\vspace{2em}図プレースホルダー：小規模コーパスからのLLMによる特徴抽出（状態／戦略／遷移）\\\texttt{fig2\_feature\_extraction} をここに挿入してください。\vspace{2em}}}
  \caption{LLMが小規模高品質コーパスの各ターンを，会話状態・応答戦略・それによって生じる状態遷移として構造化する。}
  \Description{小規模高品質コーパスがLLMによって処理され，会話状態・応答戦略・状態遷移という3つのラベル付き出力が得られることを示す図。}
  \label{fig:features}
\end{figure}

\subsection{ステップ2：ベイズ的対話モデルの構築}
\label{sec:bayesmodel}
ステップ1で得られた(状態，戦略，次状態)の3つ組は，離散状態空間$S=\{s_1,\dots,s_K\}$と離散戦略（観測ラベル）空間$O=\{o_1,\dots,o_L\}$上のベイズ的対話モデルの構築に用いられる（図\ref{fig:bayesmodel}）。このモデルは2つの要素からなる。\emph{状態遷移分布}$T(s'\mid s)$は，目的スタイルのコーパス内で，状態$s$から状態$s'$へ移行する可能性を推定する。小規模コーパス中で頻繁に見られる遷移（例えば，気持ちの開示から気持ちの整理へ，典型的には共感的な戦略を介して移行するもの）には高い遷移確率が与えられ，稀にしか観測されない遷移（例えば，気持ちを十分に探索する前に一気に解決策へ飛ぶもの）には低い確率が与えられる。\emph{観測尤度}$O(o\mid s')$は，会話が状態$s'$にあるときに，ある応答戦略$o$が観測される可能性を推定する。$T$と$O$は合わせて，小規模コーパスの特徴的な会話ダイナミクスを，表層的なテキストのレベルではなく状態と戦略のレベルで要約したコンパクトな確率モデルを構成しており，これにより小規模コーパス自体には含まれない発話に対してもモデルを再利用してスコアリングできる。

\begin{figure}[t]
  \centering
  %% \includegraphics[width=\linewidth]{figures/fig3_bayesian_model}
  \fbox{\parbox{0.9\linewidth}{\centering\vspace{2em}図プレースホルダー：状態遷移グラフ（太線＝目的コーパスで頻出する遷移，細線＝稀な遷移）\\\texttt{fig3\_bayesian\_model} をここに挿入してください。\vspace{2em}}}
  \caption{ベイズ的対話モデルは，会話状態をノード，応答戦略を状態間の確率的な遷移として表現する。太線は目的スタイルのコーパスにおいて頻出する遷移（$T(s'\mid s)$が高い）を，細線は稀な遷移を示す。}
  \Description{気持ちの開示，気持ちの整理，解決策の検討などの会話状態ノードが，遷移確率を表す太さの異なる辺で結ばれた有向グラフ。}
  \label{fig:bayesmodel}
\end{figure}

\subsection{ステップ3：ベイズ更新による候補応答のスコアリング}
\label{sec:scoring}
ベイズ的対話モデルが得られると，BASiSは与えられた会話文脈に対する候補応答$r$を，単一の最尤状態に決め打ちするのではなく，対話履歴を状態上の信念分布として扱うことでスコアリングする（図\ref{fig:scoring}）。前ターンまでの状態上の信念分布を$b_{t-1}(s)$とする。BASiSはまず，この信念を遷移モデルを通じて伝播させることで，次状態の予測分布を計算する：
\begin{equation}
  \hat{b}_t(s') = \sum_{s \in S} b_{t-1}(s)\, T(s' \mid s).
  \label{eq:predict}
\end{equation}
続いて，LLM（ステップ1と同一のラベリングスキームを用いる）が候補応答$r$を観測戦略ラベル$o\in O$に分類する。BASiSはこの観測を用いて，予測分布をベイズ則に基づいて更新する：
\begin{equation}
  b_t(s') \;=\; \frac{O(o \mid s')\, \hat{b}_t(s')}{\displaystyle\sum_{s'' \in S} O(o \mid s'')\, \hat{b}_t(s'')}.
  \label{eq:update}
\end{equation}
最後に，目的スタイルにとって望ましいとみなされる状態群$D\subset S$（例えば，来訪者の感情開示を回避せずしっかり受け止めた状態に対応するもの）を指定し，BASiSは候補応答のScoreを，その状態群への事後確率質量として定義する：
\begin{equation}
  \mathrm{Score}(r) \;=\; \sum_{s' \in D} b_t(s').
  \label{eq:score}
\end{equation}
式\eqref{eq:predict}--\eqref{eq:score}は単一の予測状態や予測戦略ではなく確率分布全体の上で動作するため，同じ文脈に対して複数の妥当な目的スタイル戦略と整合する候補応答は，たまたま1つの戦略にのみ一致する応答と比べて不利に扱われることがない。ベイズ更新は，観測と整合するすべての戦略へ自然に確率質量を分配する。このスコアリング手続きは，小規模コーパスに含まれるか否かにかかわらず任意の候補発話に適用できるため，はるかに小規模な参照モデルを用いて大規模汎用コーパス全体をスコアリングすることが可能になる。

\begin{figure}[t]
  \centering
  %% \includegraphics[width=\linewidth]{figures/fig4_scoring}
  \fbox{\parbox{0.9\linewidth}{\centering\vspace{2em}図プレースホルダー：信念伝播，候補応答のLLMによる戦略ラベリング，ベイズ更新，2つの候補例に対するScore計算\\\texttt{fig4\_scoring} をここに挿入してください。\vspace{2em}}}
  \caption{候補応答のスコアリング：前ターンまでの信念分布が遷移モデルによって予測分布へ伝播され，LLMが候補応答の戦略をラベリングし，予測分布がベイズ則で更新され，望ましい状態群への事後確率質量がScoreとなる。}
  \Description{同一プロンプトに対する2つの候補応答がスコアリングされる例。一方は望ましい状態群に事後信念質量の大半が割り当てられる高スコア応答，他方はそれ以外の状態に質量が集中する低スコア応答である。}
  \label{fig:scoring}
\end{figure}

\subsection{ステップ4：データ選別とLoRAベースDPO学習}
BASiSは，ステップ3のスコアリング手続きを，大規模汎用コーパス中の各応答候補に対して，その直前の文脈を条件として適用し，得られたScoreを用いて学習データを選別する（図\ref{fig:selection}）。プール内で複数の応答候補が存在するプロンプトについては，最もスコアの高い応答を\emph{chosen}応答，より低いスコアの応答を\emph{rejected}応答として扱い，Direct Preference Optimizationが要求する意味での選好ペアを構成する\cite{rafailov2023dpo}。これらのペアは，続いてLoRAベースのDPO\cite{hu2021lora}によってローカルLLMのファインチューニングに用いられる。LoRAアダプタは，事前学習済みの重みを凍結したまま，chosen応答をrejected応答よりも相対的に選好するようにモデルを学習する。本研究では意図的に適応量を小さく抑えている。ベースモデルの一般的な対話能力と世界知識は貴重であり損なうべきではないため，学習は1エポック，低い学習率，小さなLoRAランクを用い，積極的な再学習よりも精密なスタイル適応を優先する。学習ハイパーパラメータの詳細は第\ref{sec:setup}節に示す。

\begin{figure}[t]
  \centering
  %% \includegraphics[width=\linewidth]{figures/fig5_selection_dpo}
  \fbox{\parbox{0.9\linewidth}{\centering\vspace{2em}図プレースホルダー：大規模コーパスからのScoreに基づくデータ選別，chosen/rejectedペア構築，ローカルLLMのLoRA/DPO学習\\\texttt{fig5\_selection\_dpo} をここに挿入してください。\vspace{2em}}}
  \caption{同一プロンプトに対して高Score応答がchosen例として，低Score応答がrejected例として保持され，選好ペアを構成する。これがローカルLLMのLoRAベースDPOファインチューニングを駆動する。}
  \Description{大規模汎用コーパスがScoreによってchosen集合とrejected集合にフィルタリングされ，選好ペアへと結合され，LoRAとDPOによってローカルLLMのファインチューニングに用いられることを示す図。}
  \label{fig:selection}
\end{figure}


\section{実験設定}
\label{sec:setup}

\subsection{リサーチクエスチョン}
本予備実験は3つのリサーチクエスチョンに取り組む。\textbf{RQ1：} BASiSで選別したデータでのファインチューニングは，ファインチューニングを行わないベースモデルと比較して，目的スタイルに適合した応答の頻度を増加させるか。\textbf{RQ2：} この改善は，スタイルに基づく選別そのものに起因するのか，それとも単に同一プールから\emph{何らかの}データを用いてファインチューニングを行ったことに起因するのか――すなわち，BASiSは同一パイプラインでランダムにデータを選別した場合を上回るか。\textbf{RQ3：} この改善は，参照コーパス固有の評価軸だけでなく，より広範なNLP分野の評価実践から構成した評価軸のもとでも成立するか，またその改善は統計的に有意か。さらに2つの問い――目的特化領域を横断した汎用性，および人間の判断との一致度――は，現在進行中の実験の対象であり，第\ref{sec:limitations}節でその計画とプレースホルダーの結果表とともに述べる。

\subsection{コーパス}
\label{sec:corpora}
第\ref{sec:results}節で報告する予備実験では，ESConv\cite{liu2021esconv}を小規模高品質なスタイル参照元（1{,}254件の対話，ステップ1の特徴抽出の元データ）とし，DailyDialog\cite{li2017dailydialog}を大規模汎用な選別対象プール（ステップ3でスコアリングされる53{,}810件の応答候補）として用いた。本論文全体を通じて強調しているように，BASiSは目的特化対話全般を対象としており，ESConvはその一事例として用いている。ESConvとDailyDialogの組を用いたのは，両者が明確で既に検証済みの小規模・大規模コーパスペアであり，既存研究との比較が可能であるためである\cite{wen2025guidedtod,wan2025emodynamix}。この組み合わせを超えてBASiSが汎化するかを検証するため，現在，さらに2つの小規模・大規模コーパスペアを用いた追加実験を進めている。1つは，MathDial\cite{macina2023mathdial}（数学チュータリング対話）とMediTOD\cite{saley2024meditod}（医療問診対話）を組み合わせ，感情支援以外の2つの目的特化的な対話レジスタにまたがる小規模スタイル参照元とする設定である。もう1つは，DailyDialogに代えて，実際のユーザとLLMの対話からなるはるかに大規模で雑音の多いプールであるWildChat-1M\cite{zhao2024wildchat}を大規模コーパスとして用いる設定である。これらの設定の結果は本稿執筆時点では未取得であり，第\ref{sec:results}節では進行中として報告する。

\subsection{モデルと学習設定}
すべての条件でローカルLLMとしてQwen3.5-27Bをファインチューニングした。\textbf{Base}条件は，未修正の事前学習済みモデルである。\textbf{BASiS}条件は，式\eqref{eq:score}のScoreによってDailyDialogから選別された2{,}500件の選好ペアを用いて，LoRAベースDPOでファインチューニングしたモデルである。\textbf{Random}条件は，同一の学習手続きを用いるが，Score選別ではなくDailyDialogからランダムに選んだ2{,}500件の候補から構成した選好ペアで学習したモデルであり，学習データの量と形式を固定した上でスタイルに基づく選別のシグナルのみを取り除くことで，RQ2に対応して選別自体の寄与を分離する。DPO学習は学習率$5\times10^{-6}$，$\beta=0.1$を用い，LoRAはランク$r=8$，$\alpha=16$，dropout $0.05$で1エポック適用した。

\subsection{評価プロトコル}
\label{sec:evalprotocol}
すべての条件を，モデルの正体を伏せたブラインド条件かつ温度$0$のOracle LLM（GPT-5.4）によって，2つの相補的なルーブリックのもとで評価した。

\textbf{ESConv参照ルーブリック。} ESConvデータセット論文自体で用いられている基準\cite{liu2021esconv}に従い，Oracleは各応答を，感情の受容，助言のタイミング，支援戦略の適切さ，温かさ，会話の進行という軸に沿って評価し，そこから3つの要約指標を導出する。すなわち，感情の受容・助言のタイミング・戦略の適切さを中心とする\emph{ESConv core score}，これにトーンと会話の進行を加えた\emph{weighted overall}，そして評価対象の100件中，Oracleがどちらの条件の応答をよりESConvらしいと判断したかの割合を表す\emph{win rate}（引き分けは別に記録）である。このルーブリックは，Base対BASiS（RQ1）およびBASiS対Random（RQ2）の対応比較に用いる。

\textbf{トップカンファレンス評価実践に基づく5軸ルーブリック。} ESConv固有の基準を超えた一般性を検証するため（RQ3），Base・BASiS・Randomの3条件すべてに対し，第\ref{sec:relwork}節でレビューした評価実践から構成した5つの軸で追加評価を行った。\emph{Style Strength}は，応答がESConvらしい感情支援スタイルをどれだけ強く示しているかを表す全体的な強度\cite{mir2019styletransfer}である。\emph{ESConv Tone Similarity}は，共感的・受容的・非断定的な支援者トーンへの近さを表す\cite{liu2021esconv}。\emph{Supporter Role Consistency}は，会話履歴全体を通じて支援者としての役割・態度が維持されているかを表す\cite{zhang2018persona}。\emph{Non-directive Support Style}は，助言に急ぐのではなく，受容・整理・探索を優先しているかを表す\cite{deng2023mixedinitiative}。\emph{Premature Advice Avoidance}は，感情開示に対して，共感・受容を挟まずに即座に助言へ飛んでいないかを表す\cite{liu2021esconv,deng2023mixedinitiative}。各軸は100件のプロンプトについて1〜10点で評価され，同一の100プロンプトと各プロンプトにつき3条件分の応答（各条件1件）が，単一のブラインドパスで評価されることで，プロンプト内で対応のある比較となる。3条件間の全体的な差はFriedman検定によって各軸ごとに評価し，全体検定が有意であった場合には，各対応比較（BASiS対Base，BASiS対Random）について対応ありpermutation testを行い，複数比較に対してHolm補正を適用した。全体の効果量としてKendallの$W$を，各対応比較の効果量として対応差の標準化効果量を報告する。

\textbf{人手評価（計画中）。} Oracleによる判定は，いかに注意深くプロンプトを設計しても，人間の判断の代替とはなり得ない。特に，温かさや役割の一貫性のような人間にとって意味のある性質を測定することを意図したルーブリックにおいてはなおさらである。したがって，上記2つのルーブリックを踏襲した人手評価を，同一のプロンプト・応答集合を用いて実施するべく設計を進めている。そのプロトコルと現状は第\ref{sec:limitations}節で述べる。


\section{実験結果}
\label{sec:results}

\subsection{RQ1：BaseとBASiSの比較}
表\ref{tab:base_vs_basis}は，ESConv参照ルーブリックによるBaseモデルとBASiSファインチューニング済みモデルの比較を示す。BASiSはBaseに対し，3つの要約指標すべてで改善を示した。ESConv core scoreは6.950ポイント増加し，weighted overallは4.850ポイント増加し，評価対象100件中65\%でOracleはBASiSの応答をBaseの応答よりESConvらしいと判断した（Base優位は29\%，引き分け6\%）。これらの結果は，本予備実験の設定内において，RQ1に対する肯定的な答えを支持する。すなわち，BASiSで選別したデータでのファインチューニングは，未修正のベースモデルと比較してESConvらしい応答の頻度を増加させる。

\begin{table}[t]
  \caption{RQ1：ESConv参照ルーブリックによるBase対BASiSの比較（評価対象100件）。}
  \label{tab:base_vs_basis}
  \begin{tabular}{lccc}
    \toprule
    指標 & Base & BASiS & 差分 \\
    \midrule
    ESConv core score  & 83.576 & \textbf{90.525} & $+6.950$ \\
    Weighted overall    & 84.112 & \textbf{88.962} & $+4.850$ \\
    Win rate             & 29\%   & \textbf{65\%}   & (引き分け 6\%) \\
    \bottomrule
  \end{tabular}
\end{table}

\subsection{RQ2：BASiSとRandomの比較}
表\ref{tab:basis_vs_random}は，同一のルーブリックをBASiSとRandom選別ベースラインに適用した結果を示す。Randomは同量のDailyDialogデータで学習されているが，スタイルに基づく選別は行っていない。BASiSは3つの指標すべてでRandomを上回った。ESConv core scoreは7.011ポイント高く，weighted overallは4.770ポイント高く，OracleはプロンプトのうちBASiSの応答を62\%で選好した（Random優位31\%，引き分け7\%）。BaseとRandomは，学習データの選び方以外は同一の手続きでファインチューニングされているため，この比較はファインチューニング自体の効果からスタイルに基づく選別の効果を分離するものであり，RQ2に対して肯定的な答えを支持する。すなわち，観測された改善はBASiSの選別基準そのものに起因するものであり，単にプール内の追加データで適応を行ったことに起因するものではない。

\begin{table}[t]
  \caption{RQ2：ESConv参照ルーブリックによるBASiS対Random選別の比較（評価対象100件）。}
  \label{tab:basis_vs_random}
  \begin{tabular}{lccc}
    \toprule
    指標 & BASiS & Random & 差分 \\
    \midrule
    ESConv core score  & \textbf{90.163} & 83.152 & $+7.011$ \\
    Weighted overall    & \textbf{88.616} & 83.845 & $+4.770$ \\
    Win rate             & \textbf{62\%}   & 31\%   & (引き分け 7\%) \\
    \bottomrule
  \end{tabular}
\end{table}

\subsection{RQ3：トップカンファレンス評価実践に基づく5軸ルーブリック}
表\ref{tab:five_axis}は，Friedman検定の全体$p$値およびHolm補正付き対応比較の結果とともに，3条件すべてに対する5軸ルーブリックの結果を示す。BASiSはすべての軸で最高の平均スコアを達成した。全体のFriedman検定は5軸すべてで$p<.05$で有意であり，BASiS対Baseの対応比較は5軸すべてで（Holm補正後も）有意であった。BASiS対Randomの対応比較は5軸中4軸で有意であった。Non-directive Support Styleについては，BASiSとRandomの差（8.34対7.99）はHolm補正後に有意水準に達しなかった（$p=.0502$）。これは，慣例的な$\alpha=.05$の閾値において，偶然による差である可能性を排除できない唯一の軸である。全体の効果量（Kendallの$W$）および各対応比較の標準化効果量は，統計分析の最終版とあわせて確定作業中であり，表\ref{tab:five_axis}では保留として示す。上記の有意性の判定は，この分析から既に得られている$p$値に基づいて報告している。

\begin{table*}[t]
  \caption{RQ3：トップカンファレンスの評価実践から構成した5軸ルーブリック，3条件すべて（評価対象100件，Oracleによるブラインド1〜10点評価，同一プロンプトでの対応比較，Friedman検定とHolm補正付き対応ありpermutation test）。}
  \label{tab:five_axis}
  \footnotesize
  \begin{tabular}{lccccccc}
    \toprule
    軸 & Base & BASiS & Random & Friedman $p$ & BASiS対Base & BASiS対Random & 効果量（$W$／標準化差） \\
    \midrule
    Style Strength               & 8.26 & \textbf{8.65} & 8.21 & $<.001$  & 有意，Holm $p<.001$   & 有意，Holm $p<.001$   & 確定中 \\
    ESConv Tone Similarity       & 8.21 & \textbf{8.52} & 8.13 & $<.001$  & 有意，Holm $p<.001$   & 有意，Holm $p<.001$   & 確定中 \\
    Supporter Role Consistency   & 8.34 & \textbf{8.60} & 8.41 & $.0012$  & 有意，Holm $p=.0021$  & 有意，Holm $p=.0282$  & 確定中 \\
    Non-directive Support Style  & 7.84 & \textbf{8.34} & 7.99 & $.0055$  & 有意，Holm $p=.0051$  & 非有意，Holm $p=.0502$ & 確定中 \\
    Premature Advice Avoidance   & 8.57 & \textbf{9.44} & 8.92 & $<.001$  & 有意，Holm $p<.001$   & 有意，Holm $p=.0098$  & 確定中 \\
    \bottomrule
  \end{tabular}
  \\[2pt]
  {\footnotesize \emph{注：}「有意」「非有意」は，$\alpha=.05$におけるHolm補正付き対応ありpermutation testの結果を表す。Kendallの$W$および各対応比較の標準化効果量は確定作業中であり，カメラレディ版で挿入する。}
\end{table*}

\subsection{追加コーパスペア（進行中）}
表\ref{tab:pending}は，第\ref{sec:corpora}節で述べた追加の小規模・大規模コーパスペア――小規模スタイル参照元としてのMathDial+MediTOD，大規模選別対象プールとしてのWildChat-1M――について，第\ref{sec:evalprotocol}--\ref{sec:results}節と同一のプロトコル（ESConv参照指標は，チュータリング／医療問診設定に適した領域固有の指標に置き換える）で報告する予定の表形式を示す。これらの実験は現在進行中であり，計画している分析の完全性を査読者が判断できるよう，ここでは表の構造のみを示す。カメラレディまでに値を埋める予定である。

\begin{table}[t]
  \caption{追加コーパスペア：結果は現在進行中のため未取得（投稿時点）。}
  \label{tab:pending}
  \begin{tabular}{lccc}
    \toprule
    設定 & Base & BASiS & Random \\
    \midrule
    MathDial+MediTOD $\to$ 選別対象プール（詳細確定中） & --- & --- & --- \\
    WildChat-1M 選別対象プール（小規模コーパス確定中） & --- & --- & --- \\
    \bottomrule
  \end{tabular}
\end{table}

\subsection{人手評価（計画中）}
第\ref{sec:evalprotocol}節のESConv参照ルーブリックおよび5軸ルーブリックを踏襲した人手評価研究を計画しており，上記のOracleに基づく知見を裏付けることを目的としている。詳細は第\ref{sec:limitations}節で述べる。結果は本稿執筆時点ではまだ得られていない。


\section{考察}
\label{sec:discussion}

第\ref{sec:results}節の結果のパターンは，BASiSが，参照コーパスの中核的なスタイル的コミットメントに最も直接結びついた側面を強化することに特に有効であることを示唆している。BaseおよびRandomの双方に対する最大の改善は，Premature Advice AvoidanceとStyle Strengthで見られ，感情の受容・助言のタイミング・戦略の適切さから構成されるESConv core score自体も，表\ref{tab:base_vs_basis}の比較の中で最大の単一指標の差を示している。これは，第\ref{sec:scoring}節のスコアリング関数の設計と整合的である。Scoreの計算に用いる望ましい状態群$D$は，来訪者の開示を解決策へ移る前にしっかり受容・探索した状態に対応するよう選ばれているため，BASiSで選別したデータでのファインチューニングがまさにこれらの振る舞いを最も強く改善することは驚くべきことではない。

Non-directive Support Styleは，BASiSが生の平均値では最高であったにもかかわらず，BASiS対Randomの差がHolm補正後に有意水準に達しなかった唯一の軸である。この結果と整合する説明は2つあり，互いに排他的ではない。第一に，DailyDialogは一般対話コーパスであり，明示的に取引的なやり取りを除けば，指示的で助言優先の応答はもともと比較的少ないため，ランダムに選別した部分集合であっても偶然に中程度の非指示性スコアを示しやすく，達成可能な差が圧縮されている可能性がある。第二に，$n=100$件という評価規模では，観測された程度の真の効果（8.34対7.99）は，Holm補正付きpermutation testが$\alpha=.05$で検出できる限界に近く，人手評価で計画しているようなより大規模な評価サンプルによって，これが真の天井効果か検出力不足によるものかが明らかになる可能性がある。

分析の過程でより本質的な限界として浮かび上がったのは，受容的な振る舞いと会話の前進とのあいだのトレードオフである。表\ref{tab:five_axis}のトップカンファレンス5軸ルーブリックには含まれていないが，本研究のESConv参照評価では，各応答が（例えば適切なタイミングでの明確化質問を通じて）会話をどれだけ前に進められているかも追加で評価している。この観点では，BASiSはRandomよりも低いスコアを示した（64.30対84.29）。妥当な解釈としては，BASiSはその構造上，受容的で承認的な状態に留まる応答を報酬するため，受容と会話を前進させる積極的な質問とを組み合わせた応答を，十分に報酬しきれていない可能性がある。結果として，会話の勢いを感情的な同調と引き換えにしてしまっているとも考えられる。これは看過できない限界である。感情を際限なく受容し続けるだけで来訪者を前に進める手助けをしない感情支援対話は，実際にはESConvのスタイルを忠実に再現したものとは言えない。なぜなら，ESConvの支援者も最終的には会話を前進させるからである。この結果は，単一の望ましい状態群から構築された単一のスカラーScoreでは，複数の，時に競合するサブ的振る舞いのバランスから成り立つ目的スタイルを十分に捉えきれない可能性を示している。本研究では，これをBASiSを改善する上で最も重要な具体的方向性と位置づけている。すなわち，受容志向の状態に加えて前進志向の状態も含むよう望ましい状態群の定義をより広く・多要素的にするか，あるいは単一のスカラーScoreから明示的な多目的選別基準へと移行するかのいずれかである。本研究では，このトレードオフを，第\ref{sec:setup}節で述べた追加コーパスペア実験における設計上の目標の1つと位置づけている。チュータリングや医療問診のような，感情支援よりもさらに会話の前進が本質的に重要となり得る領域は，このトレードオフに対する自然なストレステストを提供する。

より広く見れば，これらの知見は，目的特化的な会話スタイルが一枚岩ではないことを示唆している。単一の目的コーパス内であっても，スタイルは複数のサブ要素から構成されており，1つの要素に向けて調整された選別手法によって，それらが不均等に――本研究の場合は部分的には互いを犠牲にする形で――改善され得る。これは，BASiSあるいは類似の手法を適用しようとする者にとって実践的な含意を持つ。すなわち，式\eqref{eq:score}における望ましい状態群$D$の定義は，一度固定して領域を問わず使い回すべき細部ではなく，対象領域ごとに見直すべき設計上の選択であり，理想的にはその領域の目的にとってどのサブ的振る舞いが最も重要かという知見に基づいて決定されるべきである。

第\ref{sec:results}節の結果の妥当性に対する脅威として，さらに2点を指摘しておく。第一に，現時点で報告しているすべての比較は単一のOracle判定モデルに依拠している。ブラインド評価，固定の温度，適切な統計的検定を用いることで，見せかけの結果のリスクは低減しているが，LLM判定器には（例えば，より長い応答やより強調的な応答を好むといった）系統的バイアスが存在することが知られており，計画中の人手評価は，まさにこのリスクに対処することを目的としている。第二に，本稿で報告している定量的な結果は，単一の領域（感情支援）と単一のベースモデル（Qwen3.5-27B）に基づくものである。BASiSの汎用性について確信を持って主張するためには，追加のコーパスペア実験，そして将来的には追加のベースモデルによる検証が必要である。


\section{限界と進行中の取り組み}
\label{sec:limitations}

以下では，本研究が抱える限界を率直に報告するとともに，それぞれに対して現在講じている，あるいは講じる予定の具体的な対応を述べる。

\textbf{単一の評価者。} 第\ref{sec:results}節のすべての結果は，単一のLLM Oracle判定器に基づいている。第\ref{sec:evalprotocol}節の2つのルーブリックの双方を踏襲した人手評価研究を設計中であり，同一の100プロンプトとプロンプトごとの3応答を，同一のブラインド対応比較デザインのもとで人間の評価者が同一の軸に沿って評価する。これにより，Oracleに基づく結果と人手に基づく結果を直接比較でき，Oracleのみの評価では得られない評価者間一致率などの統計量も得られる。この研究は現在進行中であり，その結果，および結果に応じて生じ得る主張の修正は，本論文の次バージョンで報告する予定である。

\textbf{単一の完全に実施された領域。} 第\ref{sec:results}節の定量的な結果は，すべてESConv／DailyDialogの組み合わせから得られたものである。第\ref{sec:corpora}節で述べたとおり，チュータリングと医療問診のレジスタにまたがる小規模スタイル参照元としてMathDialとMediTODを組み合わせ，DailyDialogよりもはるかに大規模かつ雑音の多い大規模選別対象プールとしてWildChat-1Mを用いる追加実験を進めている。これらの実験により，BASiSの改善が感情支援を超えて汎化するか，また第\ref{sec:discussion}節で議論した受容／前進のトレードオフがこの領域固有のものか，より一般的に生じるものかを検証する。

\textbf{特徴抽出の誤りの伝播。} BASiSのベイズ的対話モデルは，全面的にLLMが生成した状態・戦略・遷移ラベルから構築されている（第\ref{sec:method}節ステップ1）。このラベリングにおける誤りや不整合は，遷移分布・観測分布に直接伝播し，そこから下流のすべてのScoreへと波及する。ラベリングの信頼性（例えば，複数回のラベリングパス間の一致度，あるいはラベリングLLMと人間のアノテータとの一致度）について，本稿執筆時点では体系的な監査を行っていない。このような監査は，計画中の人手評価研究の自然な補完であり，現在進行中の取り組みの一部である。

\textbf{単一のベースモデル。} 現時点でのすべてのファインチューニング実験はQwen3.5-27Bを用いている。BASiSの改善が他のモデルファミリーやパラメータ規模でも成立するかは未検証であり，今後の課題である。

\textbf{受容／前進のトレードオフ。} 第\ref{sec:discussion}節で議論したとおり，BASiSの現行の望ましい状態群の定義は，会話の前進に対して一定の代償を払いながら，受容的で承認的な振る舞いを改善している。本研究では，これを本研究が提起した最も重要な未解決の方法論的問いと位置づけており，望ましい状態群の定義を拡張・多要素化するか，あるいは明示的な多目的スコアリング基準へ移行するかのいずれかによって対処する予定であり，上述の追加領域実験において評価する。


\section*{倫理およびプライバシーに関する声明}
本研究は，感情支援のような感情的に機微な目的，および進行中の拡張においては医療問診に関連する会話スタイルへ向けて言語モデルを適応させるものである。第\ref{sec:results}節で報告した実験では，公開されている既に匿名化された研究用コーパス（ESConv，DailyDialog，MathDial，MediTOD，WildChat-1M）のみを使用し，人間の参加者から新たにデータを収集することは行っていない。第\ref{sec:limitations}節で述べた計画中の人手評価は，人間を対象とするデータの収集に先立ち，適切な倫理審査およびインフォームド・コンセントの手続きのもとで実施する。支援的な会話スタイルを模倣するようファインチューニングされたシステムは，専門的なメンタルヘルスケアや医療の代替とはならず，本研究はBASiSを実運用可能な支援ツールや臨床ツールとして提示するものではない。本手法に基づく将来的な実運用には，領域に応じた安全性評価，危機的状況にあるユーザのためのエスカレーション経路，システムがAIであることの明確な開示が必要であり，これらはいずれも本データ選別研究の対象範囲外である。本手法固有のデュアルユースリスクは，対話データでLLMをファインチューニングすることに一般的に伴うリスクを超えて大きくはないと考えているが，このようなスタイル選別技術は，原理的には支援的なスタイルではなく操作的・欺瞞的な会話スタイルを模倣するために応用され得ることには留意すべきである。今後の研究および実運用の意思決定において，この可能性が明示的に考慮されることを望む。


\section*{GenAI利用開示}
本研究では，大規模言語モデルを，執筆支援のみならず研究方法論の中核的な構成要素として利用した。BASiS自体が，小規模コーパスからの会話特徴抽出（第\ref{sec:method}節ステップ1），スコアリング時の候補応答の戦略ラベリング（第\ref{sec:scoring}節），および自動評価のOracle判定器（第\ref{sec:evalprotocol}節）にLLMを利用しており，これらの利用はいずれも本手法に本質的なものであり，上記の該当箇所で詳細に記述している。\emph{[著者への注記：以下の記述はテンプレートであり，ACM IUI 2027のGenAI利用開示要件に従って，投稿前に著者自身が内容を確認・完成させる必要があります。]} 本原稿の一部の文章は，著者の直接の監督のもとで大規模言語モデルの支援を受けて起草された。技術的な主張，実験設計上の判断，および報告されている結果はすべて，著者が作成・検証し，その責任を負うものである。LLMの支援を受けて作成されたコードや分析スクリプトは，使用前に著者によってレビュー・検証されている。表\ref{tab:base_vs_basis}〜\ref{tab:five_axis}に報告されている数値結果は，上記および第\ref{sec:evalprotocol}節で述べたOracle判定器としての役割を除き，LLMによって生成または改変されたものではない。


\section{結論}
本研究では，LLMによって構造化されたベイズ的対話モデルを用いて，大規模汎用コーパスからスタイル適合的なデータをスコアリング・選別し，LoRAベースDPOファインチューニングに用いることで，ローカルLLMの会話スタイルを小規模高品質コーパスのそれへ適応させる手法BASiSを提案した。ガイドラインに基づく応答制御とは異なり，BASiSは人手によるルール設計を必要としない。単一戦略の応答予測とは異なり，妥当な戦略の自然な多様性を単一ラベルではなく確率分布として表現する。ESConv参照の感情支援スタイルを対象とした予備実験では，BASiSで選別したデータが，領域固有のルーブリックとより広範な評価実践から構成した5軸ルーブリックの双方において，ファインチューニングを行わないベースモデルおよびランダム選別ベースラインの双方に対してスタイル適合度を改善し，複数比較に対する適切な補正のもとでもその大半が統計的に有意であった。また，これらの改善に留保を付ける受容／前進のトレードオフも明らかになり，本研究ではこれを今後の手法改良における中心的な未解決課題として位置づけている。現在進行中の取り組みでは，BASiSを感情支援以外の目的特化領域および人手評価へと拡張しており，目的特化対話システムにおけるスタイルに基づくデータ選別の汎用性を確信をもって確立するために，これらはいずれも必要な検証であると考えている。


%% ------------------------------------------------------------------
%% ACM IUI 2027の投稿規定により，GenAI利用開示および参考文献はワード数に含まれません。
%% ------------------------------------------------------------------
\bibliographystyle{ACM-Reference-Format}
\bibliography{basis_ja}

\end{document}
